\section{はじめに} \label{sec1}

近年，Embodied AI~\cite{duan2022}と呼ばれる，身体性を持つエージェントが現実世界や仮想空間で学習を行い，多様な物体が存在する複雑な環境でさまざまなタスクを遂行しようとする研究が行われている．例えば，「エージェントが人間のある指示に従って家事（例：掃除，洗濯）をする」，「エージェントが環境の中にある特定の物を見つけて案内を行う」といったタスクが挙げられる．この研究分野では，自然言語の指示に従って，エージェントがタスクを実行するためのタスクプランニング（行動計画）をするために，大規模言語モデル（Large Language Models, LLMs）を活用することが注目されている~\cite{huang22,rana2023sayplan,song2023llm}．これらの研究により，LLMは日常生活でタスクを実行する上で重要な常識知識を有することが実証されてきた．

LLMを用いて，環境の状況（オブジェクト（物体）の状態や位置，物体間の位置関係）に応じてエージェントの行動を計画する研究も進められている~\cite{lin2023grounded}．
このような行動計画をLLMによって実現するためには，LLMに環境の状況を正確に認識させることが課題となる．しかし，環境が広範・複雑であるほど情報量が多くなるため，すべての情報をLLMに認識させることは非現実的である．そのため，必要な情報を絞り込むことが課題となる．

また，既存研究で用いられるデータセット~\cite{Puig_2018_CVPR,ALFRED20}は，環境内の状況を完全に把握していなくても，常識的な知識を用いることで達成可能なタスクを含んでいる．つまり，これらのタスクでは，物体の状態や識別番号，物体間の関係といった詳細な知識を必要としない場合が多い．例えば，``Turn on the tv''のようなタスクを遂行するためには，テレビまで近づき，テレビの電源をオンにする行動計画を生成することが一般的に正解と考えられる．しかし，このタスクは単一の物体（テレビ）を対象としており，また，その物体の電源がオフであることが前提となっている．そのため，LLMは，電源がオフであるという暗黙の常識知識に基づき，前述の行動計画を生成することができる．一方，``Turn on all light switches''のようなタスクでは，家庭内に存在するライトの数や，それぞれのライトの状態（オンまたはオフ）と識別番号を認識していなければ，適切な行動計画を生成することが困難である．このように，既存研究では，家庭環境を完全に把握しなければ解決が難しいタスク設定はされていない．

本研究では，環境内の物体の状態や物体間の関係などを含む詳細な情報を家庭環境知識としてLLMに認識をさせ，家庭環境知識に従ってタスクを実行するための適切な行動計画を生成することを目的としている．提案手法では，LLMと家庭環境知識を用いて，タスクを遂行するための適切な行動計画を逐次的に生成する．また，行動計画の構成要素であるアクションステップを生成する過程で，生成したアクションを実行するための前提条件を満たしているかをLLMを用いて検証する．この前提条件を満たしている場合には，アクションを実行する．この前提条件を満たしていない場合には，LLMを用いて抽出した家庭環境知識をもとに，再度アクションステップを生成する．ここで，行動計画とは，タスクを遂行するために必要な一連のアクションを順序立てた計画を指し，アクションステップとは，行動計画を構成する各個別の行動を指す．

提案手法の評価を行うために，本研究では，家庭内での活動を三次元仮想空間内でシミュレートするマルチエージェントプラットフォームVirtualHome~\cite{Puig_2018_CVPR}を用いて，家庭環境の把握が不可欠な家庭内タスクのデータセットを作成する．具体的には，物体の状態に着目した「状態変化タスク」と，物体間の関係に着目した「配置タスク」の2種類を対象とする．これらのタスクでは，コンピュータビション分野の研究のように，エージェント視点の視覚情報を処理して環境を認識する手法~\cite{song2023llm,kim2023context,nakajima2024combining}ではなく，シミュレータの内部から取得したデータを家庭環境知識として用いることを前提とする．また，本研究では，家庭環境知識をLLMに認識させることで，家庭環境知識に基づいた適切な行動計画を生成する手法を提案する．

本論文の貢献は以下の通りである．
\begin{itemize}
    \item 家庭環境を把握していなければ解決することが困難である家庭内タスクの設計およびデータセットの構築．
    \item LLMが家庭環境を認識し，その理解に基づいてタスク実行のための行動計画を生成する手法の提案．
    \item 行動計画を生成する過程で，生成したアクションを実行するための前提条件の検証および再生成を行う新しいプロセスの導入．
\end{itemize}

本論文の構成は次の通りである．第~\ref{sec2}では，本研究の関連研究について述べる．第~\ref{sec3}では，タスク設定とデータセット構築について述べる．第~\ref{sec4}では，本研究の提案手法とその構成について述べる．第~\ref{sec5}では，評価実験，結果，考察について述べる．最後に第~\ref{sec6}では，本研究の結論と今後の課題について述べる．

なお，本論文は，青山らの研究成果~\cite{青山2024,Aoyama2025}に，追加実験を行った上で改訂したものである．
また，提案手法のソースコードおよび構築したデータセットは，GitHubにて公開している~\footnote{\url{https://github.com/ke-lab-it-agu/Precheck-LLM-Planner}}．